% =====================================================================
% RELATÓRIO 9 — POO com Ruby
% Capítulo para ser incorporado à Apostila S01 (OO + PF), após o Relatório 8.
% =====================================================================

% ---------------------------------------------------------------------
% Definições usadas neste capítulo (podem ser movidas para o preâmbulo).
% - ruby: habilita a linguagem Ruby no listings e corrige a acentuação
%   dentro dos códigos (mesma solução dos estilos csharp e python).
% - azulRelatorio e exercicio: já são criados no Relatório 7. As linhas
%   abaixo só os recriam se este capítulo for compilado sem ele.
% ---------------------------------------------------------------------
\lstdefinestyle{ruby}{
    language=Ruby,
    morekeywords={include,raise},
    literate=
        {á}{{\'a}}1 {à}{{\`a}}1 {ã}{{\~a}}1 {â}{{\^a}}1
        {é}{{\'e}}1 {ê}{{\^e}}1 {è}{{\`e}}1 {í}{{\'\i}}1
        {ó}{{\'o}}1 {õ}{{\~o}}1 {ô}{{\^o}}1 {ú}{{\'u}}1
        {ç}{{\c{c}}}1 {Á}{{\'A}}1 {É}{{\'E}}1 {Ç}{{\c{C}}}1
        {Ã}{{\~A}}1 {Â}{{\^A}}1 {Ê}{{\^E}}1 {Õ}{{\~O}}1
        {—}{{\textemdash}}1 {Í}{{\'I}}1 {Ó}{{\'O}}1 {Ú}{{\'U}}1
}

\definecolor{azulRelatorio}{rgb}{0.1,0.3,0.8}

\ifcsname exercicio\endcsname\else
\newtcolorbox{exercicio}[1]{
    colback=white,
    colframe=azulRelatorio,
    coltitle=white,
    fonttitle=\normalfont,
    boxrule=0.8pt,
    arc=1.5mm,
    title={#1}
}
\fi

\chapter{Relatório 9: Programação Orientada a Objetos em Ruby}

\section{Fundamentos Teóricos e Pilares da POO}

Nos Relatórios 6 a 8, os pilares da POO foram apresentados em C++, C\# e Python. O Python mostrou
que boa parte das palavras-chave (\texttt{private}, \texttt{virtual}, \texttt{override}) pode ser
trocada por convenções. O \textbf{Ruby} vai além: nele, \textbf{tudo é objeto}, inclusive
números, textos e as próprias classes, e a sintaxe foi desenhada para ser lida quase como uma
frase em inglês. Os blocos não usam chaves nem dependem da indentação: eles terminam com a
palavra \texttt{end}. Os objetos são criados com \texttt{.new}, e os atributos começam com
\texttt{@} e \textbf{nunca} podem ser acessados diretamente de fora do objeto: a própria classe
decide o que expor, criando métodos. Assim como no C\# e no Python, a memória é liberada
automaticamente pelo \textbf{Garbage Collector}.

A tabela abaixo relaciona as estruturas já vistas em Python com as equivalentes em Ruby:

\begin{center}
\begin{tabular}{|l|l|l|}
\hline
\textbf{Conceito} & \textbf{Python (Relatório 8)} & \textbf{Ruby (Relatório 9)} \\
\hline
Construtor & \texttt{def \_\_init\_\_(self, nome)} & \texttt{def initialize(nome)} \\
Atributo do objeto & \texttt{self.nome} & \texttt{@nome} \\
Instanciação & \texttt{Item(nome)} & \texttt{Item.new(nome)} \\
Atributo privado & prefixo \texttt{\_\_} (convenção) & toda variável \texttt{@} \\
Getter simples & \texttt{def get\_nome(self)} & \texttt{attr\_reader :nome} \\
Setter & \texttt{arma.set\_dano(20)} & \texttt{arma.dano = 20} \\
Herança & \texttt{class Arma(Item)} & \texttt{class Arma < Item} \\
Construtor do pai & \texttt{super().\_\_init\_\_(nome)} & \texttt{super(nome)} \\
Método do pai & \texttt{super().usar()} & \texttt{super} \\
Interface & \texttt{ABC} + \texttt{@abstractmethod} & \texttt{module} + \texttt{include} \\
Coleção & \texttt{list} e \texttt{.append()} & \texttt{Array} e \texttt{<{}<} \\
Repetição & \texttt{for x in lista:} & \texttt{lista.each do |x| ... end} \\
Fim de bloco & indentação & \texttt{end} \\
\hline
\end{tabular}
\end{center}

Os quatro pilares continuam valendo, agora ilustrados com o universo do anime
\textit{Gachiakuta}. Nele, a sociedade se divide entre quem vive nas alturas e quem é jogado no
\textbf{Abismo}, um lugar cheio de lixo. Lá vivem os \textbf{Catadores}, que dão vida a objetos
descartados, chamados \textbf{Utensílios}, e se organizam em grupos de \textbf{Zeladores}:

\begin{itemize}
    \item \textbf{Abstração}: Cada objeto guarda apenas o que importa para a história. Um
    \texttt{Habitante} do Abismo precisa de nome, origem e força.

    \item \textbf{Encapsulamento}: A força fica no atributo interno \texttt{@\_forca}, que só
    pode ser lido e alterado pelos métodos \texttt{forca} e \texttt{forca=}, e o setter recusa
    valores inválidos. Já o nome e a origem podem ser lidos de fora, mas não alterados.

    \item \textbf{Herança}: \texttt{Catador} reaproveita a estrutura de \texttt{Habitante}
    (relação ``\textbf{É UM}'': um Catador é um Habitante).

    \item \textbf{Polimorfismo}: O grupo de Zeladores percorre os seus membros chamando
    \texttt{atacar}, e cada objeto executa a sua própria versão do método.
\end{itemize}

\subsection{Módulos: Contratos e Mixins}

A novidade deste relatório é a forma como o Ruby trata as interfaces. Ele não tem a
palavra-chave \texttt{interface}, como o C\#, nem classes abstratas nativas, como o \texttt{ABC}
do Python. Em seu lugar, usa \textbf{módulos}: blocos de métodos que qualquer classe pode
\textbf{incluir} com \texttt{include}. Um módulo pode ter dois papéis:

\begin{itemize}
    \item \textbf{Contrato}: o módulo declara um método que só lança um
    \texttt{NotImplementedError}, e a classe que inclui o módulo deve escrever a sua própria
    versão do método. No código de referência, \texttt{CapacidadeEspirito} exige que o
    \texttt{Catador} implemente \texttt{invocar\_espirito}.

    \item \textbf{Mixin}: o módulo já traz o método pronto, e toda classe que o incluir ganha
    esse comportamento sem precisar reescrevê-lo. No código de referência,
    \texttt{CapacidadeCombate} dá o método \texttt{entrar\_em\_combate} ao \texttt{Catador}.
\end{itemize}

Uma classe pode incluir vários módulos ao mesmo tempo. Essa é a alternativa do Ruby à herança
múltipla vista no Python: uma classe tem um único pai, mas pode somar os comportamentos de
quantos módulos precisar.

Há uma diferença importante em relação ao Relatório 8. No Python, o \texttt{ABC} impede que um
objeto seja criado se faltar algum método abstrato. O Ruby \textbf{não confere nada
antecipadamente}: o objeto é criado normalmente, e o \texttt{NotImplementedError} só aparece no
momento em que o método que falta é chamado. Por isso, em Ruby, o que importa é se o objeto
\textbf{sabe responder} a um método, e não de qual classe ele é. Essa ideia é chamada de
\textit{duck typing} (``se anda como um pato e grasna como um pato, então é um pato''), e o
método \texttt{respond\_to?} permite fazer essa pergunta ao objeto.

\subsection{Relacionamentos entre Classes}

A composição e a agregação, vistas nos relatórios anteriores, continuam valendo (relação
``\textbf{TEM UM}''):

\begin{itemize}
    \item \textbf{Composição (relação forte)}: A parte é criada \textbf{dentro} do objeto dono e
    ``morre'' junto com ele. Exemplo: o Utensílio nasce dentro do \texttt{initialize} do próprio
    Catador.

    \item \textbf{Agregação (relação fraca)}: A parte é criada \textbf{fora} e o dono apenas
    guarda uma referência a ela. Exemplo: os Catadores são criados antes e só depois são
    entregues ao grupo de Zeladores, e continuam existindo se o grupo deixar de existir.
\end{itemize}

\section{Documentação Sintática das Estruturas de POO}

Esta seção detalha todas as palavras-chave e construções sintáticas utilizadas na implementação
de POO em Ruby. Os exercícios deste relatório devem ser resolvidos \textbf{apenas} com as
estruturas documentadas aqui.

\subsection{Estrutura do Programa e Saída no Console}

O Ruby não tem uma função \texttt{main}: o arquivo é executado de cima para baixo, e o código
solto depois das classes funciona como programa principal. Comentários começam com \texttt{\#}.
Blocos (classes, métodos, \texttt{if}, laços) terminam com \texttt{end}. O \texttt{puts} imprime
um valor e pula uma linha. Dentro de aspas duplas, \texttt{\#\{...\}} insere um valor no texto
(\textbf{interpolação}), como a f-string do Python.

\begin{lstlisting}[style=ruby]
# comentários começam com #
heroi = "Rudo"
puts "O herói se chama #{heroi}."   # interpolação com #{}
puts "Origem: " + "Aoiro"           # + junta dois textos
puts "\n--- Fim ---"                # \n = quebra de linha
\end{lstlisting}

\begin{itemize}
    \item \textbf{Variáveis}: Não têm tipo declarado. Basta atribuir um valor, como em
    \texttt{heroi = ...}.

    \item \textbf{Parênteses opcionais}: Em uma chamada de método, os parênteses podem ser
    omitidos. \texttt{puts x} e \texttt{puts(x)} fazem a mesma coisa.

    \item \texttt{+}: Entre dois textos, junta os dois em um texto só.
\end{itemize}

\subsection{Declaração de Classe, \texttt{initialize} e Atributos com \texttt{@}}

Uma classe é declarada com \texttt{class} e termina com \texttt{end}. O construtor é o método
\texttt{initialize}, executado automaticamente quando o objeto é criado com \texttt{.new}. Os
atributos do objeto são as variáveis que começam com \texttt{@} (\textbf{variáveis de
instância}). Elas não precisam ser declaradas antes: basta atribuir um valor.

\begin{lstlisting}[style=ruby]
class Utensilio

  # roda automaticamente quando o objeto é criado com .new
  def initialize(nome, tipo)
    @nome = nome   # @ indica um atributo do próprio objeto
    @tipo = tipo
  end

  # a última expressão do método é o seu valor de retorno
  def mostrar_detalhes
    "Utensilio #{@nome} do tipo #{@tipo}."
  end
end

luvas = Utensilio.new("Luvas", "Ofensiva")  # instanciação
puts luvas.mostrar_detalhes
\end{lstlisting}

\begin{itemize}
    \item \texttt{def ... end}: Declara um método. Um método sem parâmetros dispensa os
    parênteses, tanto na declaração quanto na chamada (\texttt{luvas.mostrar\_detalhes}).

    \item \texttt{initialize}: O construtor, equivalente ao \texttt{\_\_init\_\_} do Python.
    Diferente do Python, ele não recebe o \texttt{self} como parâmetro.

    \item \texttt{@nome}: Variável de instância. Cada objeto tem a sua própria cópia, que existe
    enquanto o objeto existir.

    \item \texttt{.new}: Cria o objeto e chama o \texttt{initialize} com os argumentos
    recebidos.

    \item \textbf{Retorno implícito}: O método retorna o valor da sua última expressão, sem
    precisar de \texttt{return}. A palavra \texttt{return} também existe e pode ser usada para
    sair antes do fim do método.

    \item \texttt{self}: Dentro de um método, referencia o próprio objeto (como o \texttt{self}
    do Python e o \texttt{this} do C\#), mas não aparece nos parâmetros. A expressão
    \texttt{self.class} (ou \texttt{objeto.class}) retorna a classe do objeto, e dentro de um
    texto ela aparece como o nome da classe, como em \texttt{\#\{self.class\}}.
\end{itemize}

\subsection{Encapsulamento: \texttt{attr\_reader}, \texttt{attr\_writer} e
\texttt{attr\_accessor}}

Em Ruby, toda variável de instância é \textbf{inacessível de fora do objeto}: não existe sintaxe
para ler um atributo com \texttt{@} de fora da classe. O nível de acesso é decidido pelos
\textbf{métodos} que a classe oferece. Para não ser preciso escrever um getter para cada
atributo, o Ruby oferece atalhos que criam esses métodos automaticamente:

\begin{lstlisting}[style=ruby]
class Habitante
  attr_reader :nome, :origem   # cria os métodos nome e origem

  def initialize(nome, origem, forca_inicial)
    @nome = nome
    @origem = origem
    @_forca = forca_inicial     # interno: não ganhou attr_reader
  end
end

barado = Habitante.new("Barado", "Superficie", 800)
puts barado.nome         # funciona: attr_reader criou o método nome
barado.nome = "Outro"    # NoMethodError: o método nome= não existe
\end{lstlisting}

\begin{itemize}
    \item \texttt{attr\_reader :nome}: Cria apenas o getter \texttt{nome}. De fora, o atributo
    pode ser lido, mas não alterado.

    \item \texttt{attr\_writer :nome}: Cria apenas o setter \texttt{nome=}.

    \item \texttt{attr\_accessor :nome}: Cria os dois. No código de referência, é usado no
    \texttt{utensilio} do \texttt{Catador}.

    \item \texttt{:nome}: Um \textbf{símbolo}, isto é, um nome fixo precedido de dois-pontos. Os
    símbolos são usados para se referir a métodos e atributos pelo nome.

    \item \texttt{@\_forca}: O \texttt{\_} não muda nada para o Ruby, porque toda variável com
    \texttt{@} já é inacessível de fora. Ele é só uma convenção, usada no código de referência,
    para destacar um atributo interno, que não ganhou \texttt{attr\_reader} e só é acessado
    pelos métodos da própria classe.

    \item \texttt{NoMethodError}: Erro lançado quando se chama um método que o objeto não
    possui, como \texttt{nome=} em uma classe que só tem \texttt{attr\_reader :nome}.
\end{itemize}

\subsection{Getters e Setters Manuais}

Quando o atributo precisa de validação, o getter e o setter são escritos à mão. O setter é um
método cujo nome termina com \texttt{=}; por isso, ele é chamado com a mesma sintaxe de uma
atribuição:

\begin{lstlisting}[style=ruby]
def forca              # getter: devolve o atributo interno
  @_forca
end

def forca=(novo_valor) # setter: o nome termina com =
  if novo_valor.is_a?(Numeric) && novo_valor.positive?
    @_forca = novo_valor
  end
end

rudo.forca = 1100      # chama o setter forca=
puts rudo.forca        # chama o getter forca
\end{lstlisting}

No código de referência, um valor inválido é simplesmente ignorado. Para avisar o usuário, basta
acrescentar um \texttt{else}. O exemplo abaixo é complementar e não faz parte do código de
referência:

\begin{lstlisting}[style=ruby]
def forca=(novo_valor)
  if novo_valor > 0
    @_forca = novo_valor
  else
    puts "A força deve ser positiva!"
  end
end
\end{lstlisting}

\begin{itemize}
    \item \texttt{if}, \texttt{elsif}, \texttt{else} e \texttt{end}: Estruturas de decisão.

    \item \textbf{Comparações}: \texttt{==}, \texttt{!=}, \texttt{>}, \texttt{<}, \texttt{>=} e
    \texttt{<=}. Condições podem ser combinadas com \texttt{\&\&} (e), \texttt{||} (ou) e
    \texttt{!} (não), como em \texttt{x >= 1 \&\& x <= 10}.

    \item \texttt{is\_a?(Classe)}: Retorna \texttt{true} se o objeto for da classe informada ou
    de uma filha dela (o equivalente ao \texttt{isinstance} do Python). \texttt{Numeric} é a
    classe de todos os números.

    \item \texttt{positive?}: Retorna \texttt{true} se o número for maior que zero. Por
    convenção, os métodos terminados em \texttt{?} respondem sim ou não (\texttt{true} ou
    \texttt{false}).

    \item \textbf{Setter dentro da própria classe}: Deve ser chamado como
    \texttt{self.forca = valor}. Sem o \texttt{self.}, o Ruby entende que
    \texttt{forca = valor} cria uma variável local, e o setter não é executado.
\end{itemize}

\subsection{Herança e \texttt{super(...)} no \texttt{initialize}}

A classe pai é indicada com \texttt{<} depois do nome da classe filha. Para inicializar a parte
herdada, a filha chama \texttt{super(...)} dentro do seu \texttt{initialize}, repassando os
parâmetros de que o pai precisa (o equivalente ao \texttt{super().\_\_init\_\_(...)} do Python e
ao \texttt{: base(...)} do C\#):

\begin{lstlisting}[style=ruby]
class Catador < Habitante   # Catador É UM Habitante

  def initialize(nome, origem, forca_inicial, nome_utensilio,
                 tipo_utensilio)
    # nome, origem e forca_inicial vão para o initialize de Habitante
    super(nome, origem, forca_inicial)
    @utensilio = Utensilio.new(nome_utensilio, tipo_utensilio)
  end
end
\end{lstlisting}

Se a filha \textbf{não} define um \texttt{initialize}, ela herda automaticamente o do pai e é
criada com os mesmos parâmetros. O exemplo abaixo é complementar e não faz parte do código de
referência:

\begin{lstlisting}[style=ruby]
class Andarilho < Habitante   # sem initialize próprio
  def atacar
    puts "#{nome} arremessa uma pedra!"
  end
end

# usa o initialize herdado de Habitante
andarilho = Andarilho.new("Andarilho", "Abismo", 500)
\end{lstlisting}

\begin{itemize}
    \item \texttt{class Filha < Pai}: Declara a herança (relação ``É UM'').

    \item \textbf{O que é herdado}: A filha herda os métodos do pai, inclusive os criados por
    \texttt{attr\_reader}. Por isso, o \texttt{Andarilho} pode usar \texttt{nome} diretamente.

    \item \textbf{Variáveis com} \texttt{@}: Pertencem ao objeto, e não à classe que as criou.
    Diferente do \texttt{\_\_} do Python, a filha enxerga as variáveis criadas no
    \texttt{initialize} do pai.
\end{itemize}

\subsection{Sobrescrita e \texttt{super} em Métodos}

Em Ruby não existem \texttt{virtual} nem \texttt{override}: qualquer método pode ser
sobrescrito. Basta a classe filha definir um método com o \textbf{mesmo nome} do método do pai:

\begin{lstlisting}[style=ruby]
# Na Classe Base (Habitante):
def atacar
  puts "#{nome} desfere um golpe basico!"
end

# Na Classe Derivada (Catador):
def atacar
  puts "#{nome} ataca usando o utensilio #{@utensilio.nome}!"
  puts "  [Tipo: #{@utensilio.tipo}]"
end
\end{lstlisting}

Dentro do método sobrescrito, a palavra \texttt{super} executa a versão do pai (o equivalente ao
\texttt{super().metodo()} do Python e ao \texttt{base.Metodo()} do C\#). Se a versão do pai
\textbf{retorna} um texto, a filha pode reaproveitá-lo e apenas acrescentar o que é dela. O
exemplo abaixo é complementar e não faz parte do código de referência:

\begin{lstlisting}[style=ruby]
# Na Classe Base (Habitante): retorna um texto
def apresentacao
  "#{nome}, vindo de #{origem}"
end

# Na Classe Derivada (Catador): reaproveita o texto do pai
def apresentacao
  super + " | Utensilio: #{@utensilio.nome}"
end
\end{lstlisting}

\begin{itemize}
    \item \texttt{super} (sem parênteses): Chama o método de mesmo nome no pai, repassando os
    mesmos argumentos que o método atual recebeu.

    \item \texttt{super()}: Chama o método do pai sem nenhum argumento.

    \item \texttt{super(a, b)}: Chama o método do pai com os argumentos informados. É a forma
    usada no \texttt{initialize} do \texttt{Catador}.
\end{itemize}

\subsection{Módulos como Contrato: \texttt{module}, \texttt{include} e
\texttt{NotImplementedError}}

Um módulo é declarado com \texttt{module} e incluído em uma classe com \texttt{include}. Para
funcionar como contrato, o método do módulo apenas lança um \texttt{NotImplementedError}. A
classe que inclui o módulo escreve a sua própria versão do método, que tem prioridade sobre a do
módulo:

\begin{lstlisting}[style=ruby]
module CapacidadeEspirito
  def invocar_espirito
    raise NotImplementedError,
          "#{self.class} deve implementar invocar_espirito"
  end
end

class Catador < Habitante
  include CapacidadeEspirito   # "assina" o contrato

  def invocar_espirito         # cumpre o contrato
    puts "#{nome} invoca o espirito do utensilio! Poder: #{forca * 5}"
  end
end
\end{lstlisting}

\begin{itemize}
    \item \texttt{module Nome ... end}: Declara o módulo. O nome começa com letra maiúscula,
    como o de uma classe, mas um módulo não pode ser instanciado com \texttt{.new}.

    \item \texttt{include Modulo}: Acrescenta os métodos do módulo à classe. Os objetos da
    classe passam a responder a esses métodos.

    \item \texttt{raise Erro, mensagem}: Lança um erro e interrompe o programa, exibindo a
    mensagem. A mensagem é opcional: \texttt{raise NotImplementedError} sozinho também funciona.

    \item \texttt{NotImplementedError}: Erro que indica um método ainda não implementado.

    \item \textbf{Prioridade}: Quando a classe e o módulo têm um método com o mesmo nome, vale a
    versão da classe. Por isso, o \texttt{invocar\_espirito} do \texttt{Catador} substitui o do
    módulo.
\end{itemize}

O Ruby não verifica o contrato antes da hora. O mesmo \texttt{raise NotImplementedError} pode
ser usado diretamente em um método de uma classe base, para indicar que as filhas devem
sobrescrevê-lo, mas a classe base continua podendo ser instanciada. O exemplo abaixo é
complementar e não faz parte do código de referência:

\begin{lstlisting}[style=ruby]
class Criatura
  def rugir
    raise NotImplementedError, "Cada criatura ruge do seu jeito"
  end
end

mancha = Criatura.new   # funciona: o objeto é criado normalmente
mancha.rugir            # NotImplementedError: o erro só aparece aqui
\end{lstlisting}

\subsection{Módulos com Implementação Pronta (Mixins)}

Um módulo também pode trazer métodos completos. Toda classe que o incluir ganha esse
comportamento pronto, e uma mesma classe pode incluir quantos módulos precisar:

\begin{lstlisting}[style=ruby]
module CapacidadeCombate
  def entrar_em_combate
    puts "#{self.class} entra em postura de combate!"
  end
end

class Catador < Habitante
  include CapacidadeEspirito
  include CapacidadeCombate    # uma classe, vários módulos
end
\end{lstlisting}

Dentro de um método de módulo, \texttt{self} é o objeto que está executando o método, ou seja,
um objeto da classe que incluiu o módulo. Por isso, o módulo pode chamar os métodos desse objeto,
como um getter criado por \texttt{attr\_reader}, desde que a classe que o incluir tenha esse
método. O exemplo abaixo é complementar e não faz parte do código de referência:

\begin{lstlisting}[style=ruby]
module Rastreavel
  def emitir_sinal
    # nome vem do attr_reader da classe que incluir o módulo
    puts "Sinal de #{nome} detectado no Abismo!"
  end
end
\end{lstlisting}

\subsection{Composição e Agregação}

A diferença entre os dois relacionamentos está em \textbf{onde a parte é criada}:

\begin{lstlisting}[style=ruby]
# COMPOSIÇÃO: o Utensilio é criado DENTRO do initialize do Catador
@utensilio = Utensilio.new(nome_utensilio, tipo_utensilio)

# AGREGAÇÃO: os Catadores são criados FORA e só recebidos pelo grupo
grupo = GrupoDeZeladores.new([rudo, enjin])
\end{lstlisting}

\subsection{Arrays: Criação, \texttt{<{}<} e \texttt{.each}}

O \texttt{Array} é a coleção de tamanho dinâmico do Ruby, equivalente à \texttt{list} do Python.
Como a tipagem é dinâmica, um mesmo array pode guardar objetos de classes diferentes, o que
possibilita o polimorfismo ao percorrê-lo:

\begin{lstlisting}[style=ruby]
catadores = []           # array vazio
catadores << rudo        # adiciona no final
catadores.push(enjin)    # mesma coisa que <<
puts catadores.length    # quantidade de elementos

catadores.each do |c|    # percorre elemento por elemento
  c.atacar
end

catadores.each { |c| c.atacar }   # a mesma coisa, em uma linha
\end{lstlisting}

\begin{itemize}
    \item \texttt{[]}: Cria um array vazio. \texttt{[a, b]} cria um array já preenchido.

    \item \texttt{array <{}< objeto} e \texttt{array.push(objeto)}: Inserem um objeto no final do
    array.

    \item \texttt{array.length} (ou \texttt{array.size}): Retorna a quantidade de elementos.

    \item \texttt{.each}: Executa o \textbf{bloco} uma vez para cada elemento. O bloco pode ser
    escrito entre \texttt{do} e \texttt{end} ou entre chaves, e \texttt{|c|} declara a variável
    que recebe o elemento da vez.

    \item \texttt{def initialize(catadores = [])}: Um parâmetro pode ter \textbf{valor padrão},
    usado quando o argumento não é informado.

    \item \texttt{puts array}: Imprime cada elemento do array em uma linha.
\end{itemize}

\subsection{Transformar e Filtrar: \texttt{.map}, \texttt{.select} e \texttt{respond\_to?}}

Além do \texttt{.each}, os arrays têm métodos que percorrem os elementos e devolvem um
\textbf{novo} array, sem alterar o original. As duas primeiras linhas de código do exemplo
abaixo são complementares e não fazem parte do código de referência:

\begin{lstlisting}[style=ruby]
# .map: novo array com o resultado do bloco para cada elemento
nomes = catadores.map { |c| c.nome }

# .select: novo array só com os elementos aprovados pelo bloco
fortes = catadores.select { |c| c.forca > 1000 }

# introspecção: seleciona só quem possui o método entrar_em_combate
combatentes = @catadores.select { |c| c.respond_to?(:entrar_em_combate) }
combatentes.each(&:entrar_em_combate)
\end{lstlisting}

\begin{itemize}
    \item \texttt{.map}: Devolve um novo array, do mesmo tamanho do original, com o valor da
    última expressão do bloco para cada elemento.

    \item \texttt{.select}: Devolve um novo array só com os elementos para os quais o bloco
    resultou em \texttt{true}.

    \item \texttt{objeto.respond\_to?(:metodo)}: Retorna \texttt{true} se o objeto possui o
    método informado, inclusive quando ele veio de um módulo. É a forma de perguntar ao objeto
    o que ele sabe fazer (\textbf{introspecção}).

    \item \texttt{each(\&:metodo)}: Atalho para \texttt{each \{ |c| c.metodo \}}, que chama o
    método informado em cada elemento.
\end{itemize}

\section{Código de Referência Completo}

Abaixo está a implementação de referência utilizada para exemplificação do paradigma POO em
Ruby:

\begin{lstlisting}[style=ruby, caption={Exemplo de Aplicação de POO em Ruby}]
# ------- Contexto: Gachiakuta ---------
# No mundo de Gachiakuta, a sociedade é dividida entre quem vive nas
# alturas e quem é jogado no Abismo — um lugar repleto de lixo e
# criaturas chamadas Manchas.
# Os Catadores são habitantes do Abismo que aprenderam a dar vida a
# objetos descartados, chamados Utensílios, usando o poder do espírito.
# Organizados em grupos de Zeladores, eles partem em missões para
# sobreviver e recuperar o que foi perdido nas profundezas.


# MÓDULOS (funcionam como interfaces em outras linguagens)
# definem um comportamento que qualquer classe pode "incluir"
# a classe que incluir este módulo deve implementar invocar_espirito;
# o Ruby não confere isso antes: o erro só aparece quando o método
# for chamado em um objeto que não o implementou
module CapacidadeEspirito
  def invocar_espirito
    raise NotImplementedError,
          "#{self.class} deve implementar invocar_espirito"
  end
end

# MÓDULO DE COMBATE
# outro módulo, mas esse já traz a implementação pronta
# qualquer classe que incluir CapacidadeCombate ganha o
# entrar_em_combate() de graça
module CapacidadeCombate
  def entrar_em_combate
    puts "#{self.class} entra em postura de combate!"
  end
end


# COMPOSIÇÃO
# Utensilio não existe sozinho fora do Catador
# ele nasce e morre junto com ele
class Utensilio
  attr_reader :nome, :tipo

  # Construtor
  def initialize(nome, tipo)
    @nome = nome
    @tipo = tipo
  end

  def mostrar_detalhes
    "Utensilio #{@nome} do tipo #{@tipo}."
  end
end


# CLASSE BASE
# molde principal de todo habitante do Abismo
# as classes filhas vão herdar daqui
class Habitante

  # ENCAPSULAMENTO
  # em Ruby, toda variável de instância (as que começam com @) já é
  # inacessível de fora da classe: só dá para lê-la ou alterá-la por
  # meio de métodos
  # attr_reader cria apenas o getter para nome e origem: de fora, eles
  # podem ser lidos, mas não alterados
  # o _ em @_forca não muda nada para o Ruby: é só uma convenção para
  # indicar um atributo interno, que não ganhou attr_reader
  # como precisamos de validação na força (precisa ser um número
  # positivo), criamos getter e setter manuais
  attr_reader :nome, :origem

  def initialize(nome, origem, forca_inicial)
    @nome = nome
    @origem = origem
    @_forca = forca_inicial
  end

  # GETTER e SETTER manual
  # como @_forca é interno, criamos métodos pra ler e alterar ele de
  # forma controlada
  def forca
    @_forca
  end

  def forca=(novo_valor)
    # o if aqui funciona igual ao de outras linguagens:
    # só atualiza @_forca se novo_valor for um número positivo
    if novo_valor.is_a?(Numeric) && novo_valor.positive?
      @_forca = novo_valor
    end
  end

  # VIRTUAL (equivalente)
  # método que as classes filhas vão sobrescrever com o comportamento
  # específico delas
  def atacar
    puts "#{nome} desfere um golpe basico!"
  end
end


# HERANÇA + MÓDULOS
# Catador herda de Habitante e inclui os dois módulos
# por isso deve implementar invocar_espirito (contrato do módulo)
class Catador < Habitante
  include CapacidadeEspirito
  include CapacidadeCombate

  # attr_accessor é como ter attr_reader e attr_writer juntos
  # ou seja, pode ler E alterar o utensílio de fora da classe
  # usamos accessor aqui porque o Catador pode trocar de Utensílio
  attr_accessor :utensilio

  # SUPER
  # super chama o initialize do pai (Habitante) passando os dados que
  # ele precisa
  # depois inicializamos o que é específico do Catador
  def initialize(nome, origem, forca_inicial, nome_utensilio, tipo_utensilio)
    super(nome, origem, forca_inicial)

    # COMPOSIÇÃO em ação
    # Utensilio é criado aqui dentro, ele pertence ao Catador
    @utensilio = Utensilio.new(nome_utensilio, tipo_utensilio)
    puts "#{nome} surge com o #{@utensilio.nome}."
  end

  # SOBRESCRITA (override)
  # reescreve o atacar() com o comportamento específico do Catador
  def atacar
    puts "#{nome} ataca usando o utensilio #{@utensilio.nome}!"
    puts "  [Tipo: #{@utensilio.tipo}]"
  end

  # IMPLEMENTAÇÃO DO MÓDULO
  # aqui a classe cumpre o contrato do CapacidadeEspirito
  # sem isso, chamar invocar_espirito lançaria o NotImplementedError
  # definido no módulo
  def invocar_espirito
    puts "#{nome} invoca o espirito do utensilio! Poder: #{forca * 5}"
  end
end


# AGREGAÇÃO
# GrupoDeZeladores guarda referências pra Catadores que já existiam
# fora dele
# se o grupo sumir, os catadores continuam existindo
class GrupoDeZeladores
  attr_accessor :catadores

  def initialize(catadores = [])
    @catadores = catadores
  end

  def iniciar_missao(missao)
    puts "\n--- Grupo iniciando missao: #{missao} ---"

    # POLIMORFISMO EM AÇÃO
    # o mesmo .atacar() para todos, mas cada um executa o seu próprio
    @catadores.each do |c|
      c.atacar
    end
  end

  def listar_combatentes
    puts "\n--- Zeladores prontos para combate ---"

    # INTROSPECÇÃO com .select
    # .select percorre o array e retorna só os elementos que passam na
    # condição do bloco
    # o bloco é o trecho entre {} — é como um mini-método anônimo
    # executado pra cada elemento
    # |c| é a variável temporária que representa o elemento atual da
    # iteração
    # as barras || são a sintaxe do Ruby pra declarar esse parâmetro do
    # bloco, igual um (c) em outras linguagens
    # respond_to?(:entrar_em_combate) verifica se o objeto c possui
    # aquele método disponível
    combatentes = @catadores.select { |c| c.respond_to?(:entrar_em_combate) }
    combatentes.each(&:entrar_em_combate)
  end
end


# --- MAIN ---

rudo   = Catador.new("Rudo",   "Aoiro",  1000, "Luvas",        "Ofensiva")
enjin  = Catador.new("Enjin",  "Abismo", 1200, "Guarda-chuva", "Velocidade")
barado = Habitante.new("Barado", "Superficie", 800)

rudo.forca = 1100
puts "\nRudo agora tem #{rudo.forca} de forca."

puts "\n--- Habilidades ---"
rudo.invocar_espirito
rudo.entrar_em_combate
# barado.invocar_espirito <- daria NoMethodError: Habitante não inclui
# o módulo CapacidadeEspirito

puts "\nComposicao:"
puts rudo.utensilio.mostrar_detalhes

# agregação: catadores já existiam antes do grupo ser criado
grupo = GrupoDeZeladores.new([rudo, enjin])

grupo.iniciar_missao("Recuperar utensilios perdidos no Abismo")
grupo.listar_combatentes
\end{lstlisting}

\section{Exercícios Práticos do Módulo}

\begin{exercicio}{Exercício 1: Encapsulamento e Herança (O Cardápio do Afterlife)}
Em \textit{Cyberpunk 2077}, o Afterlife é o bar mais famoso de Night City, uma metrópole
futurista, e é frequentado por mercenários. O cardápio tem coquetéis comuns e drinks que
homenageiam \textbf{lendas}. Todo drink tem um nome, que nunca muda, e um preço, que o bar pode
reajustar dentro da tabela da casa.

\textbf{Glossário:}
\begin{itemize}
    \item \textbf{Lenda}: mercenário famoso de Night City, vivo ou morto, homenageado com um
    drink. No código, é só um texto.
    \item \textbf{Eurodólar}: a moeda de Night City. No código, o preço é só um número.
\end{itemize}

Implemente:

\begin{itemize}
    \item \texttt{CoquetelAfterlife}: o \texttt{initialize} recebe \texttt{nome} e
    \texttt{preco}. O \texttt{nome} pode ser lido de fora, mas não alterado (use
    \texttt{attr\_reader}). O preço fica guardado em \texttt{@\_preco}, sem
    \texttt{attr\_reader}, e o \texttt{initialize} atribui o valor recebido diretamente a ele. A
    classe tem o getter manual \texttt{preco}; o setter manual \texttt{preco=(valor)}, que só
    aceita valores de 5 a 200 (fora dessa faixa, exibe ``Preço inválido!'' e mantém o valor
    anterior); e o método \texttt{ficha}, que \textbf{retorna} (não imprime) um texto com o nome
    e o preço do drink.

    \item \texttt{DrinkDeLenda}: herda de \texttt{CoquetelAfterlife} e possui o atributo
    exclusivo \texttt{lenda} (com \texttt{attr\_reader}). O \texttt{initialize} recebe
    \texttt{nome}, \texttt{preco} e \texttt{lenda} e repassa os dois primeiros ao pai com
    \texttt{super(nome, preco)}. Sobrescreve \texttt{ficha}, reaproveitando o texto do pai com
    \texttt{super} e acrescentando a lenda homenageada.
\end{itemize}

No programa principal, crie um array com um \texttt{CoquetelAfterlife} e dois
\texttt{DrinkDeLenda}, à sua escolha, e percorra-o com \texttt{.each}, imprimindo a
\texttt{ficha} de cada drink. Depois, tente definir, com o setter, o preço de um dos drinks de
lenda como 500 e depois como 80, e imprima a ficha dele de novo. Por fim, tente alterar o nome de um dos drinks
diretamente (\texttt{objeto.nome = ...}) e observe o erro. Depois, transforme essa linha em
comentário (\texttt{\#}) para o programa voltar a rodar.
\end{exercicio}

\begin{exercicio}{Exercício 2: Abstração, Polimorfismo e Arrays (A Orquestra da Gala)}
Em \textit{Your Lie in April} (\textit{Shigatsu wa Kimi no Uso}), jovens músicos de música
clássica se apresentam em competições e concertos de gala. Em uma apresentação, cada
instrumentista tem um papel: o \textbf{solista} toca a melodia principal e o
\textbf{acompanhador} toca a base, seguindo o ritmo do solista. Todos são conduzidos pelo
\textbf{regente}.

\textbf{Glossário:}
\begin{itemize}
    \item \textbf{Solista}: músico que toca a melodia principal da peça, em destaque.
    \item \textbf{Acompanhador}: músico que toca a base da peça, seguindo o ritmo do solista.
    \item \textbf{Regente}: quem conduz a orquestra durante a apresentação (o maestro).
    \item \textbf{Peça}: a obra musical apresentada. No código, é só um texto com o nome dela.
\end{itemize}

Implemente:

\begin{itemize}
    \item \texttt{Instrumentista} (abstração): o \texttt{initialize} recebe \texttt{nome} e
    \texttt{instrumento} (os dois com \texttt{attr\_reader}). O método
    \texttt{interpretar(peca)} tem como corpo apenas um \texttt{raise NotImplementedError}, com
    uma mensagem à sua escolha.

    \item \texttt{Solista}: herda de \texttt{Instrumentista}, \textbf{sem} \texttt{initialize}
    próprio (ela usa o do pai), e sobrescreve \texttt{interpretar(peca)} para imprimir o nome, o
    instrumento e a peça, dizendo que ele conduz a melodia principal.

    \item \texttt{Acompanhador}: também herda de \texttt{Instrumentista}, \textbf{sem}
    \texttt{initialize} próprio, e sobrescreve \texttt{interpretar(peca)} para imprimir o nome,
    o instrumento e a peça, dizendo que ele acompanha o ritmo do solista.

    \item \texttt{Regente}: o \texttt{initialize} recebe \texttt{nome} e cria o array
    \texttt{@orquestra} vazio. Possui o método \texttt{convocar(instrumentista)}, que adiciona o
    instrumentista ao \texttt{@orquestra} com \texttt{<{}<}; o método \texttt{reger(peca)}, que
    imprime o nome do regente e a peça e, com \texttt{.each}, chama \texttt{interpretar(peca)}
    de cada instrumentista; e o método \texttt{escalacao}, que usa \texttt{.map} para
    \textbf{retornar} um novo array de textos, um por instrumentista, com o nome da classe dele
    (use \texttt{.class}) e o seu nome.
\end{itemize}

No programa principal, crie um \texttt{Solista}, um \texttt{Acompanhador} e um
\texttt{Regente}, à sua escolha. Convoque os dois instrumentistas, chame \texttt{reger(...)}
com uma peça à sua escolha e imprima, com \texttt{.each}, cada texto do array retornado por
\texttt{escalacao}. Por fim, crie um \texttt{Instrumentista} diretamente e, na linha seguinte,
chame \texttt{interpretar(...)} nele. Observe que o objeto é criado, mas a chamada gera um erro.
Depois, transforme a linha da chamada em comentário (\texttt{\#}) para o programa voltar a
rodar.
\end{exercicio}

\begin{exercicio}{Exercício 3: Composição e Agregação (D-Mail para o Passado)}
Em \textit{Steins;Gate}, um grupo de amigos mantém o Laboratório de Gadgets do Futuro, um
pequeno laboratório de invenções em Tóquio. A invenção mais importante do grupo é o
\textbf{PhoneWave}, um micro-ondas ligado a um celular que, por acidente, envia mensagens de
texto para o passado. O aparelho foi montado dentro do laboratório e só existe nele
(composição). Os membros do laboratório já eram pessoas comuns antes de entrar no grupo e
continuam existindo fora dele (agregação).

\textbf{Glossário:}
\begin{itemize}
    \item \textbf{PhoneWave}: o ``Microondas-Telefone'', nome que o aparelho recebe no anime.
    \item \textbf{D-Mail}: mensagem de texto enviada para o passado pelo PhoneWave. No código, é
    só um texto.
    \item \textbf{Lab Mem}: abreviação de \textit{Laboratory Member}. Cada membro do laboratório
    recebe um número de membro. No código, esse número é um inteiro.
\end{itemize}

Implemente:

\begin{itemize}
    \item \texttt{PhoneWave}: método \texttt{enviar\_dmail(mensagem)}, que imprime a mensagem e
    avisa que ela foi enviada para o passado.

    \item \texttt{MembroDoLab}: o \texttt{initialize} recebe \texttt{numero} (o número de Lab
    Mem) e \texttt{nome}, os dois com \texttt{attr\_reader}.

    \item \texttt{LaboratorioDeGadgets}: o \texttt{initialize} não recebe parâmetros, cria um
    \texttt{PhoneWave} \textbf{dentro} dele, guardado em \texttt{@phonewave}, e cria o array
    \texttt{@membros} vazio. Possui o método \texttt{admitir(membro)}, que adiciona um membro já
    criado ao \texttt{@membros} e imprime o número e o nome dele; e o método
    \texttt{experimento(mensagem)}, que imprime, com \texttt{.each}, o número e o nome de cada
    membro presente e depois envia a mensagem pelo \texttt{enviar\_dmail} do PhoneWave guardado
    em \texttt{@phonewave}.
\end{itemize}

No programa principal, crie dois membros \textbf{antes} do laboratório, à sua escolha. Depois,
crie o laboratório, admita os dois e chame \texttt{experimento(...)} com uma mensagem à sua
escolha. Explique, em um comentário no código, o que é composição e o que é agregação no seu
programa.
\end{exercicio}

\begin{exercicio}{Exercício 4: Módulos e Introspecção (Depoimentos em Baker Street)}
Sherlock Holmes é um detetive que mora na Baker Street, em Londres. Ao investigar um crime, ele
ouve todos os envolvidos, e cada um presta \textbf{depoimento} do seu jeito: a
\textbf{testemunha ocular} conta o que viu e o suspeito apresenta o seu \textbf{álibi}. No fim,
Holmes entrega à \textbf{Scotland Yard} quem pode ser detido, e só os suspeitos podem.

\textbf{Glossário:}
\begin{itemize}
    \item \textbf{Depoimento}: relato de uma pessoa sobre o crime. \textbf{Depoente} é quem
    presta o depoimento.
    \item \textbf{Testemunha ocular}: pessoa que viu o crime acontecer.
    \item \textbf{Álibi}: a justificativa do suspeito de que estava em outro lugar na hora do
    crime.
    \item \textbf{Scotland Yard}: a polícia de Londres.
\end{itemize}

Implemente:

\begin{itemize}
    \item \texttt{Depoente} (módulo-contrato): método \texttt{depor}, cujo corpo é apenas um
    \texttt{raise NotImplementedError}.

    \item \texttt{Detivel} (módulo com implementação pronta): método \texttt{ser\_detido}, que
    imprime o \texttt{nome} da pessoa e que ela foi levada pela Scotland Yard. O \texttt{nome}
    vem do \texttt{attr\_reader} da classe que incluir o módulo.

    \item \texttt{Envolvido} (classe base): o \texttt{initialize} recebe \texttt{nome} (com
    \texttt{attr\_reader}). Ela não inclui nenhum módulo.

    \item \texttt{TestemunhaOcular}: herda de \texttt{Envolvido} e inclui \textbf{apenas}
    \texttt{Depoente}. Implementa \texttt{depor} para imprimir o nome e o que a testemunha viu
    (texto à sua escolha).

    \item \texttt{SuspeitoDoCrime}: herda de \texttt{Envolvido} e inclui \texttt{Depoente}
    \textbf{e} \texttt{Detivel}. Implementa \texttt{depor} para imprimir o nome e o álibi do
    suspeito (texto à sua escolha).

    \item \texttt{Investigacao} (agregação): o \texttt{initialize} recebe \texttt{caso} (o nome
    do caso) e cria o array \texttt{@envolvidos} vazio. Possui o método
    \texttt{registrar(envolvido)}, que adiciona um envolvido já criado ao array; o método
    \texttt{ouvir\_todos}, que imprime o nome do caso e chama \texttt{depor} de cada envolvido
    com \texttt{.each} (polimorfismo); e o método \texttt{acionar\_scotland\_yard}, que usa
    \texttt{.select} com \texttt{respond\_to?(:ser\_detido)} para separar quem pode ser detido e
    chama \texttt{ser\_detido} de cada um deles.
\end{itemize}

No programa principal, crie uma testemunha ocular e dois suspeitos, à sua escolha,
\textbf{antes} da investigação. Depois, crie a investigação com um caso à sua escolha, registre
os três e chame \texttt{ouvir\_todos} e \texttt{acionar\_scotland\_yard}.
\end{exercicio}
